%\section{Tverrfaglige team og Agent-basert KI}

% Arvid, 2024-07-12, 2025-05-29

Kunstig intelligens kan spille en viktig rolle i å støtte samarbeid mellom tverrfaglige team av leger, psykologer, sykepleiere, fysioterapeuter og andre helsearbeidere for å løse diagnostiske problemer, planlegge og gjennomføre behandling. For eksempel kan et \textit{agentbasert KI-system} som AutoGen, LlamaIndex, LangChain eller CrewAI (merk: denne utviklingen går fort og nye agent-baserte KI-systemer blir stadig lansert)  brukes til å organisere og koordinere innsatsen til ulike spesialister i diagnostisering og behandling av for eksempel hjerneslag.

I diagnostiseringsfasen kan KI-agenter tildeles roller som representerer ulike spesialister, som nevrologer, radiologer og akuttmedisinere. Disse agentene kan analysere pasientdata, som MR-bilder, CT-skanninger og kliniske observasjoner, og dele sine funn med hverandre. En koordinerende agent kan sammenstille informasjonen og foreslå en diagnose basert på det samlede bildet.

For planlegging og gjennomføring av rehabilitering etter hjerneslag kan KI-systemet inkludere agenter som representerer fysioterapeuter, ergoterapeuter, logopeder og andre relevante spesialister. Disse agentene kan samarbeide om å utvikle en skreddersydd rehabiliteringsplan basert på pasientens spesifikke behov og fremgang. KI-systemet kan også overvåke pasientens fremgang over tid og foreslå justeringer i behandlingsplanen etter behov.

Andre KI-løsninger kan inkludere bruk av maskinlæring for å analysere store mengder pasientdata og identifisere mønstre som kan være nyttige i diagnostisering og behandlingsplanlegging. Naturlig språkprosessering kan brukes til å ekstrahere relevant informasjon fra medisinske journaler og forskningslitteratur, noe som kan støtte beslutningsprosessen for helsepersonell.

Ved å integrere disse KI-verktøyene i klinisk praksis, kan helsepersonell få verdifull støtte i sitt arbeid, noe som potensielt kan føre til mer presise diagnoser, mer effektive behandlingsplaner og bedre pasientutfall.\\


%Kunstig intelligens (KI) kan spille en viktig rolle i å støtte samarbeid mellom tverrfaglige team av leger, psykologer og andre helsearbeidere for å løse diagnostiske problemer og planlegge behandling. For eksempel kan et agentbasert KI-system som CrewAI brukes til å organisere og koordinere innsatsen til ulike spesialister i diagnostisering og behandling av hjerneslag.

%I diagnostiseringsfasen kan KI-agenter tildeles roller som representerer ulike spesialister, som nevrologer, radiologer og akuttmedisinere. Disse agentene kan analysere pasientdata, som MR-bilder, CT-skanninger og kliniske observasjoner, og dele sine funn med hverandre. En koordinerende agent kan sammenstille informasjonen og foreslå en diagnose basert på det samlede bildet.

%For planlegging og gjennomføring av rehabilitering etter hjerneslag kan KI-systemet inkludere agenter som representerer fysioterapeuter, ergoterapeuter, logopeder og andre relevante spesialister. Disse agentene kan samarbeide om å utvikle en skreddersydd rehabiliteringsplan basert på pasientens spesifikke behov og fremgang. KI-systemet kan også overvåke pasientens fremgang over tid og foreslå justeringer i behandlingsplanen etter behov.


%\begin{definisjonsbox}
%\textbf{Agentbasert KI} er en tilnærming til kunstig intelligens der multiple, autonome agenter samhandler for å løse komplekse problemer. Hver agent er en selvstendig enhet med spesifikke ferdigheter, kunnskap og mål. I kontekst av helsevesenet kan disse agentene representere ulike medisinske spesialister, diagnoseverktøy eller til og med pasienter.
%\end{definisjonsbox}

\begin{definisjonsbox}
\textbf{Agentbasert KI} er en tilnærming der flere autonome KI-agenter med spesialiserte roller og kompetanser samarbeider for å løse komplekse problemer. Hver agent fungerer som en selvstendig ``digital spesialist´´ som kan analysere data, ta beslutninger og kommunisere med andre agenter. I helsevesenet kan agentene representere ulike medisinske fagområder som arbeider sammen om diagnostikk, behandlingsplanlegging eller pasientoppfølging.
\end{definisjonsbox}

\vspace{2mm}

Agentbaserte KI-systemer har stort potensiale innen helsesektoren. Vi vil beskrive denne type KI i mer detalj og gi konkrete scenarier knyttet her til \texttt{CrewAI} (\small\url{https://www.crewai.com}).\\


\noindent \textbf{CrewAI} er et rammeverk for agentbasert KI som muliggjør koordinert samarbeid mellom ulike KI-agenter. Rammeverket er bygget opp av flere nøkkelkomponenter som sammen skaper et fleksibelt og effektivt system for problemløsning:

\begin{itemize}
\item \textbf{Agenter (Agents):} Hver agent representerer en spesialisert rolle med definerte ferdigheter, kunnskaper og ansvarsområder. I en medisinsk kontekst kan agenter representere ulike spesialister som nevrologer, radiologer eller fysioterapeuter. Agenter er autonome og kan ta beslutninger basert på sin spesialkompetanse.

\item \textbf{Oppgaver (Tasks):} Oppgaver definerer spesifikke arbeidsoppgaver som må utføres for å løse det overordnede problemet. Hver oppgave har klare mål, forventede resultater og kriterier for suksess. Oppgaver kan være sekvensielle eller parallelle, avhengig av arbeidsflyten.

\item \textbf{Verktøy (Tools):} Verktøy gir agentene tilgang til eksterne ressurser og funksjoner de trenger for å utføre sine oppgaver. Dette kan inkludere databaser, analysealgoritmer, søkemotorer eller spesialiserte medisinske verktøy som bildeanalyseprogrammer.

\item \textbf{Crew (Mannskap):} En crew er en samling av agenter som arbeider sammen mot et felles mål. Crewet koordinerer agentenes aktiviteter og sikrer at informasjon deles effektivt mellom medlemmene.

\item \textbf{Prosesser (Processes):} Prosesser definerer hvordan oppgaver utføres og koordineres innenfor crewet. CrewAI støtter ulike prosesstyper, inkludert sekvensielle prosesser (der oppgaver utføres i en bestemt rekkefølge) og hierarkiske prosesser (der en lederagent koordinerer andre agenter).
\end{itemize}

\noindent Styrken til CrewAI ligger i dens evne til å kombinere disse komponentene på en måte som etterligner hvordan tverrfaglige team arbeider i virkeligheten. Systemet muliggjør naturlig informasjonsdeling, delegering av oppgaver og kollektiv problemløsning, noe som gjør det særlig egnet for komplekse helserelaterte utfordringer.

\begin{eksempelbox}
\textbf{Praktisk eksempel - Hjerneslagdiagnostikk:}
I et CrewAI-system for hjerneslagdiagnostikk kan følgende komponenter samarbeide:
\begin{itemize}
\item \textit{Nevrolog-agent:} Analyserer kliniske symptomer og nevrologiske tester
\item \textit{Radiolog-agent:} Tolker MR- og CT-bilder
\item \textit{Koordinator-agent:} Sammenstiller funn og foreslår diagnose
\item \textit{Oppgaver:} Bildeanalyse, symptomvurdering, diagnosestilling
\item \textit{Verktøy:} Bildeanalyse-algoritmer (typisk: konvolusjonelle nettverk, CNNs), medisinske databaser, beslutningsstøttesystemer (typisk: LLMs med resonnering og dyp tenking)
\end{itemize}
\end{eksempelbox}

\vspace{2mm}

\noindent CrewAI skiller seg fra andre agentbaserte systemer ved sin fokus på rollebasert samarbeid og strukturert arbeidsflyt, noe som gjør det intuitivt for helsepersonell å forstå og implementere i klinisk praksis.


\begin{comment}
\noindent \textbf{CrewAI} er et rammeverk for agentbasert KI som muliggjør koordinert samarbeid mellom ulike KI-agenter. Rammeverket er bygget opp av flere nøkkelkomponenter som sammen skaper et fleksibelt og effektivt system for problemløsning:

\begin{enumerate}
\item Agenter:
\begin{nobullets}
   \item \textit{Roller}: Hver agent i CrewAI har en definert rolle som beskriver dens spesialområde og ansvar. For eksempel kan en agent ha rollen som ``Nevrologspesialist'' eller ``Dataanalytiker''.
   \item \textit{Mål}: Agenter har spesifikke mål som styrer deres handlinger og beslutninger. Et mål kan være ``Identifisere tidlige tegn på hjerneslag'' eller ``Optimalisere rehabiliteringsplanen''.
   \item \textit{Egenskaper}: Agenter kan ha ulike egenskaper som påvirker deres ytelse, som ``presisjon'', ``hastighet'' eller ``kreativitet''.
\end{nobullets}
\item Oppgaver:
\begin{nobullets}
   \item - CrewAI definerer oppgaver som konkrete problemer eller delmål som må løses. En oppgave kan være ``Analyser MR-bilde'' eller ``Generer behandlingsalternativer''.
   \item - Oppgaver kan ha prioriteter, tidsfrister og avhengigheter til andre oppgaver.
\end{nobullets}
\item Verktøy:
\begin{nobullets}
   \item - \textbf{Verktøy} i CrewAI er funksjoner eller ressurser som agentene kan bruke for å utføre oppgaver. Dette kan inkludere:
     \begin{itemize}
     \item Maskinlæringsmodeller for bildeanalyse
     \item Naturlig språkprosessering for journalgjennomgang
     \item Databaser med medisinske retningslinjer
     \item  Simuleringsverktøy for behandlingsutfall
     \end{itemize}
\end{nobullets}
\item Prosesser:
\begin{nobullets}
   \item \textit{Arbeidsflyt}: CrewAI definerer prosesser for hvordan agenter skal samarbeide og dele informasjon. Dette kan inkludere sekvensielle, parallelle eller iterative arbeidsmønstre.
   \item \textit{Beslutningsmekanismer}: Rammeverket inkluderer metoder for hvordan agenter skal komme til enighet, for eksempel gjennom avstemming, konsensus eller hierarkisk beslutningstaking.
   \item \textit{Læringsmekanismer}: Prosesser for hvordan systemet kan lære av tidligere erfaringer og forbedre seg over tid.
\end{nobullets}
\item Kommunikasjonsprotokoll:
\begin{nobullets}
   \item - CrewAI definerer standardiserte metoder for hvordan agenter utveksler informasjon og resultater.
   \item - Dette kan inkludere formater for dataoverføring, meldingsutveksling og rapportering.
\end{nobullets}
\item Overvåking og evaluering:
\begin{nobullets}
   \item - Rammeverket inkluderer mekanismer for å overvåke agentenes ytelse og systemets samlede effektivitet.
   \item - Dette muliggjør kontinuerlig forbedring og tilpasning av systemet.
\end{nobullets}
\item Menneske-KI-interaksjon:
\begin{nobullets}
   \item - CrewAI har grensesnitt for hvordan menneskelige eksperter kan interagere med og overstyre systemet ved behov.
   \item - Dette sikrer at kritiske beslutninger alltid har menneskelig tilsyn.\\
\end{nobullets}
\end{enumerate}

\noindent I en \textbf{medisisnk og helsefaglig kontekst} kan CrewAI implementeres som følger:

\begin{nobullets}
\item - Et team av KI-agenter representerer ulike medisinske spesialister (nevrolog, radiolog, fysioterapeut).
\item - Oppgaver defineres basert på pasientens tilstand og behandlingsforløp.
\item - Verktøy inkluderer spesialiserte diagnostiske algoritmer og behandlingsplanleggingsverktøy.
\item - Prosesser styrer hvordan agentene samarbeider om diagnostisering og behandlingsplanlegging.
\item - Kommunikasjonsprotokollen sikrer effektiv informasjonsdeling mellom agenter og med menneskelige helsearbeidere.
\item - Overvåkingssystemer evaluerer kontinuerlig kvaliteten på diagnoser og behandlingsresultater.
\item - Leger og annet helsepersonell kan gripe inn og justere systemets anbefalinger basert på deres kliniske vurdering.
\end{nobullets}

\noindent Denne strukturen gjør CrewAI til et kraftig og fleksibelt rammeverk for å håndtere komplekse medisinske scenarioer, hvor ulike spesialister må samarbeide tett for å gi optimal pasientbehandling.\\



\noindent I vårt \textbf{hjerneslag scenario} kan vi da introdusere:

\begin{enumerate}
\item Diagnostisk team: Et CrewAI-system for diagnostisering av hjerneslag kan best av følgende agenter:
\begin{nobullets}
    \item - Nevrologagent: Analyserer kliniske symptomer og nevrologiske undersøkelser.
   \item - Radiologagent: Tolker MR- og CT-bilder for å identifisere tegn på hjerneslag.
   \item - Laboratoriagent: Analyserer blodprøver og andre biologiske markører.
   \item - Anamnese-agent: Gjennomgår pasientens sykehistorie og risikofaktorer.
   \item - Koordinatoragent: Sammenstiller informasjon fra alle agenter og foreslår en diagnose.
\end{nobullets}
   Disse agentene kan jobbe parallelt og dele informasjon i sanntid, noe som potensielt kan føre til raskere og mer nøyaktige diagnoser.

\item Behandlingsplanlegging: For  utvikle en omfattende behandlingsplan, kan CrewAI inkludere:
\begin{nobullets}
   \item - Medikamentagent: Foreslår egnede medisiner basert på diagnose og pasientprofil.
   \item - Kirurgiagent: Vurderer behovet for og risikoen ved kirurgiske inngrep.
   \item - Rehabiliteringsagent: Utarbeider en langsiktig rehabiliteringsplan.
   \item - Ernæringsagent: Gir anbefalinger om kosthold for å støtte restitusjon.
   \item - Pasientprofil-agent: Tar hensyn til pasientens preferanser, livsstil og sosiale faktorer.
\end{nobullets}
\item Rehabiliteringsteam: Under rehabiliteringsfasen kan CrewAI koordinere:
\begin{nobullets}
   \item - Fysioterapiagent: Planlegger og overvåker fysisk terapi.
   \item - Ergoterapiagent: Fokuserer på å gjenopprette daglige funksjoner.
   \item - Logopediagent: Håndterer språk- og taleproblemer.
   \item - Progresjonsagent: Overvåker pasientens fremgang og justerer planen deretter.
   \item - Motivasjonsagent: Gir oppmuntring og støtte for å opprettholde pasientens engasjement.
\end{nobullets}
\end{enumerate}

\noindent For ytterligere å konkretisere dette scenariet innenfor rammen av \texttt{CrewAI}, kan vi illustrer med en skisse av strukturen in Python-kode:

%\begin{footnotesize}
%\begin{lstlisting}
\begin{minted}{python}
from crewai import Agent, Task, Crew, Process
from verktoy import MRAnalysator, BlodproveAnalysator, PasienthistorieAnalysator

# Definer agentene
nevrolog = Agent(
    rolle='Nevrolog',
    maal='Nøyaktig diagnostisere hjerneslag og anbefale behandling',
    bakgrunn='Erfaren nevrolog spesialisert innen hjerneslagdiagnostikk og behandling',
    detaljert=True
)

radiolog = Agent(
    rolle='Radiolog',
    maal='Analysere hjerneavbildning for å identifisere tegn på hjerneslag',
    bakgrunn='Ekspert på tolking av MR- og CT-bilder for nevrologiske tilstander',
    detaljert=True
)

laboratorieteknikker = Agent(
    rolle='Laboratorieteknikker',
    maal='Analysere blodprøver for hjerneslagrelaterte mark@rer',
    bakgrunn='Dyktig i å tolke komplekse blodprøveresultater',
    detaljert=True
)

# Definer oppgavene
analyser_mr = Task(
    beskrivelse='Analyser MR-bilder for tegn på hjerneslag',
    agent=radiolog,
    verktoy=[MRAnalysator()]
)

analyser_blodprover = Task(
    beskrivelse='Tolk blodprøveresultater for indikatorer på hjerneslag',
    agent=laboratorieteknikker,
    verktoy=[BlodproveAnalysator()]
)

gjennomga_pasienthistorie = Task(
    beskrivelse='Gjennomgå pasienthistorie for risikofaktorer for hjerneslag',
    agent=nevrolog,
    verktoy=[PasienthistorieAnalysator()]
)

diagnostiser_og_anbefal = Task(
    beskrivelse='Diagnostiser hjerneslagtype og anbefal behandling basert på 
      all tilgjengelig informasjon',
    agent=nevrolog
)

# Opprett teamet
hjerneslag_diagnose_team = Crew(
    agenter=[nevrolog, radiolog, laboratorieteknikker],
    oppgaver=[analyser_mr, analyser_blodprover, gjennomga_pasienthistorie, diagnostiser_og_anbefal],
    prosess=Process.sekvensiell
)

# Kjoer teamet
resultat = hjerneslag_diagnose_team.start()

print(resultat)
\end{minted}
%\end{lstlisting}
%\end{footnotesize}


\noindent Denne koden demonstrerer hvordan \texttt{CrewAI} kan brukes til å modellere et komplekst medisinsk scenario som hjerneslagdiagnose og behandlingsplanlegging. 
Implementasjonen inkluderer:
\begin{nobullets}
    \item (i) Definisjon av tre agenter: nevrolog, radiolog og laboratorieteknikker.
    \item (ii) Fire oppgaver: MR-analyse, blodprøveanalyse, gjennomgang av pasienthistorie, og diagnose med behandlingsanbefaling.
    \item (iii Opprettelse av et ``crew'' (team) som organiserer agentene og oppgavene.
    \item (iv) En sekvensiell prosess for å utføre oppgavene i en bestemt rekkefølge.
\end{nobullets}

For å gjøre denne implementasjonen fullt funksjonell, ville man trenge å implementere de spesifikke verktøyene (MRAnalysator, BlodproveAnalysator, PasienthistorieAnalysator) og definere mer detaljert logikk for agentenes samhandling og informasjonsdeling.\\

\end{comment}




\noindent I alle disse scenariene kan CrewAI-systemet kontinuerlig lære og tilpasse seg basert på nye data og resultater. For eksempel kan systemet oppdage mønstre i pasientresponser som menneskelige klinikere kanskje overser, og dermed bidra til å forbedre behandlingsprotokollene over tid.\\

En av de største fordelene med å bruke et verktøy som CrewAI i helsevesenet er muligheten for sømløs integrering av ulike spesialiteter og datakilder. Systemet kan raskt prosessere og analysere store mengder informasjon fra ulike kilder, inkludert elektroniske helsejournaler, vitenskapelige publikasjoner og kliniske retningslinjer, for å gi et helhetlig bilde av pasientens tilstand og optimale behandlingsmuligheter.\\

Implementeringen av slike systemer reiser imidlertid viktige etiske og praktiske spørsmål. Personvern, datasikkerhet og ansvarlighet er kritiske hensyn som må adresseres. Dessuten er det viktig å understreke at agentbaserte KI-systemer som CrewAI er ment å støtte, ikke erstatte, menneskelige helsearbeidere. Det endelige ansvaret for medisinske beslutninger bør alltid ligge hos kvalifisert helsepersonell.\\

\subsubsection*{Oppsummering}
Fremover kan vi forvente at agentbaserte KI-systemer vil spille en stadig viktigere rolle i helsevesenet, ikke bare i diagnostisering og behandling av akutt og alvorlig sykdom, men også i håndteringen av ulike komplekse medisinske, helsefaglige og helseøkonomiske tilstander og prosesser. Disse systemene har potensial til å forbedre pasientbehandling, redusere feil og øke effektiviteten i helsevesenet som helhet. Se også bokens digitale ressurs for flere eksempler og mer detaljert og oppdatert informasjon.



 \subsection*{Spørsmål og oppgaver}
\begin{enumerate}
\item Hvordan bør du og kan du inkludere en klinisk nevropsykolg i 
``hjerneslag-teamet''?
\item Drøft minst to andre scenarier du kan kjenne deg igjen i og spesifiser \textit{agenter}, \textit{verktøy}, \textit{mål} og \textit{prosesser} innen disse scenarier gjennom bruk av et agentbasert KI rammeverk.
\item Hva ser du som (i) de største muligheter, (ii) de største farer, og (iii) de største utfordringer med de nye agentbaserte KI-systemer innen medisin og helsefag?
\end{enumerate}

\subsection*{Kodeeksempler}
I kodeeksemplene som følger med avsnittet om \textit{Tverrfaglige team og Agent-basert KI} kan du selv utforske mulighetene og utfordringene. Eksempler og lenker finner du under {\small \url{https://github.com/arvidl/KI-i-helsefagene/tree/main/Agentbasert-KI}} (foreløpig addresse) i bokens digitale ressurs.\\

Det første eksempel, er ukontroversielt og viser hvordan du programatisk kan lage og finjustere et diagram i Python som fremstiller 
et \textbf{tverrfaglig team} nettverk i en helsefag-sammenheng \\ 
(se notebooken {\small \texttt{tverrfaglige-team-eksempel-0.ipynb}} som bruker \texttt{NetworkX}-biblioteket \cite{Hagberg2008}). \\

Du kan gjøre et dypdykk inn i tematikken og teknologien ved å eksperimentere med et funksjonsrikt \textbf{agentbasert KI system for diagnostikk og rehabilitering av hjerneslag}, inspirert fra tverrfaglige team slik vi kjenner fra klinikken \\ (se notebooken {\small \texttt{tverrfaglige-team-eksempel-1.ipynb}} som benytter \texttt{CrewAI} biblioteket for orkestrering av agenter, verktøy og oppgaver). \\




